\documentclass[../main.tex]{subfiles}
\begin{document}

\section{Introduction}

\label{sec:sentiment_intro}

Sentiment Analysis, also known as opinion mining, is a fundamental task in Natural Language Processing (NLP) that aims to automatically identify and extract subjective information from text, determining the emotional tone behind a piece of writing \cite{jin2019textfooler}. This computational task has become increasingly critical in our digital age, where understanding public opinion, customer feedback, and social media sentiment drives decision-making across industries ranging from marketing and finance to politics and public health \cite{morris2020textattack}. Traditional sentiment analysis systems have evolved from simple rule-based approaches to sophisticated neural network architectures that have demonstrated remarkable performance on benchmark datasets \cite{miyato2017virtual}.

Despite their impressive capabilities on benchmark datasets, recent research has exposed a concerning vulnerability in these neural sentiment analysis models: their susceptibility to adversarial attacks \cite{jin2019textfooler}\cite{li2018textbugger}. Adversarial examples in the text domain represent carefully crafted input modifications that, while appearing natural and semantically coherent to human readers, can deceive machine learning models into making incorrect predictions \cite{gao2018deepwordbug}. Unlike adversarial attacks in computer vision where pixel-level perturbations can be imperceptible, text-based adversarial examples face unique challenges due to the discrete nature of language and the need to maintain grammatical correctness and semantic meaning \cite{morris2020textattack}.

The implications of these vulnerabilities are particularly severe for sentiment analysis applications deployed in security-sensitive contexts \cite{jin2019textfooler}. Recent incidents demonstrate the practical severity of these vulnerabilities: adversarial text manipulation has contributed to market volatility when sentiment analysis systems misclassified coordinated social media campaigns, and content moderation systems have been evaded through character-level perturbations that maintain toxic intent while appearing benign to automated filters \cite{vazquez2023white}\cite{oroy2024adversarial}. In such scenarios, an adversary could potentially manipulate market sentiment or evade content filters by crafting adversarial text that misleads the underlying sentiment classification models while appearing benign to human moderators \cite{gao2018deepwordbug}.

To systematically investigate these vulnerabilities, our research focuses on three prominent adversarial attack methodologies that have demonstrated effectiveness in black-box settings \cite{morris2020textattack}. \textbf{TextFooler} \cite{jin2019textfooler} employs a synonym-based substitution strategy, identifying the most important words in the input text and replacing them with semantically similar alternatives that preserve the overall meaning while maximizing the likelihood of misclassification. This approach leverages pre-trained word embeddings and uses semantic similarity constraints to ensure that the generated adversarial examples remain coherent and natural \cite{jin2019textfooler}. \textbf{TextBugger} \cite{li2018textbugger} takes a different approach by applying character-level transformations such as character insertion, deletion, and substitution, combined with word-level modifications, making it particularly effective against real-world systems and more robust to simple preprocessing defenses \cite{li2018textbugger}. Finally, \textbf{DeepWordBug} \cite{gao2018deepwordbug} operates as a black-box attack that uses character-level perturbations guided by scoring strategies to identify critical tokens, requiring only access to the model's final predictions without knowledge of internal parameters or gradients \cite{gao2018deepwordbug}.

The choice of these three attack methods is strategic, as they collectively represent a spectrum of practical adversarial text generation techniques, including word-level and character-level perturbations, as well as heuristic-based approaches. By implementing these attacks using the \textbf{TextAttack library} \cite{morris2020textattack}, a comprehensive framework that provides standardized implementations and evaluation metrics for adversarial attacks in NLP, we ensure reproducible and fair comparisons across different attack strategies \cite{jin2019textfooler}. It is important to note that, while some of these attacks were originally proposed with white-box access in mind, our implementation using TextAttack operates in a black-box setting, requiring only access to the model's predictions, which better reflects real-world deployment constraints.

To counteract these threats, we investigate a comprehensive defense strategy combining two complementary approaches: \textbf{Virtual Adversarial Training (VAT)} \cite{miyato2017virtual} and \textbf{character-level purification}. VAT is a regularization technique that improves model robustness by generating virtual adversarial examples during training—perturbations that are designed to maximally change the model's output distribution while maintaining proximity to the original input \cite{miyato2017virtual}. Unlike traditional adversarial training that requires actual adversarial examples, VAT generates these perturbations on-the-fly using gradient information, making it computationally efficient and applicable to various neural network architectures \cite{miyato2017virtual}. To address the limitations of VAT against character-level attacks, we complement it with a lightweight character-level purification defense that operates as a preprocessing layer, targeting specific character-level transformations such as homoglyph substitutions, character insertions, deletions, and adjacent swaps commonly employed by attacks like TextBugger and DeepWordBug \cite{morris2020textattack}.

Our experimental methodology focuses on measuring the effectiveness of these attacks through accuracy degradation metrics, evaluating how significantly each attack method can reduce the performance of sentiment analysis models \cite{jin2019textfooler}. The systematic evaluation allows us to understand the relative strengths and weaknesses of different attack approaches, providing insights into the vulnerability landscape of modern sentiment analysis systems \cite{li2018textbugger}. Subsequently, we assess the defensive capabilities of our hybrid defense strategy by measuring the accuracy recovery achieved when models are trained with VAT regularization and protected by character-level purification \cite{miyato2017virtual}.

The significance of this research extends beyond academic interest, as sentiment analysis systems are increasingly deployed in real-world applications where security and reliability are paramount \cite{morris2020textattack}. Despite the widespread deployment of sentiment analysis in production systems processing millions of texts daily, systematic evaluation of adversarial robustness remains limited, with current literature focusing predominantly on transformer-based models while traditional architectures like BiLSTM-GloVe combinations remain prevalent in resource-constrained environments \cite{belmoukadam2023adversnlp}. Understanding the attack surface and developing effective defenses is crucial for maintaining trust in these systems and ensuring their safe deployment in production environments \cite{jin2019textfooler}. 

Our research makes three primary contributions: (1) systematic evaluation of adversarial vulnerabilities in BiLSTM-GloVe sentiment models across word-level, character-level, and hybrid attack strategies; (2) implementation and assessment of Virtual Adversarial Training defense effectiveness; and (3) development of a novel hybrid defense combining VAT with character-level purification, achieving up to 76.91\% attack success rate reduction against character-level attacks \cite{miyato2017virtual}\cite{morris2020textattack}.

In the subsequent sections, we will detail our experimental setup, present comprehensive results from our attack implementations, analyze the effectiveness of our hybrid defense strategy, and discuss the broader implications of our findings for the security and robustness of sentiment analysis systems in practical deployment scenarios \cite{morris2020textattack}.





\section{Methodology}

\label{sec:sentiment_methodologyy}

This section details the experimental design, implementation methodology, and evaluation framework employed to investigate adversarial vulnerabilities in sentiment analysis models and assess the effectiveness of Virtual Adversarial Training (VAT) as a defense mechanism. Our methodology encompasses the complete pipeline from model architecture selection to adversarial attack implementation and defense evaluation.

\subsection{Target Sentiment Analysis Model Architecture}

\subsubsection{Model Design and Architecture}

Our target sentiment analysis model employs a BiLSTM architecture combined with pre-trained GloVe embeddings, representing a classical yet effective approach to sentiment classification tasks \cite{pimpalkar2022mblstmglove}. The model architecture consists of the following components:

\begin{enumerate}
\item \textbf{Embedding Layer}: Utilizes pre-trained GloVe 6B.100d embeddings with a vocabulary size of 50,000 words, providing rich semantic representations learned from large-scale text corpora \cite{li2022glove}
\item \textbf{Bidirectional LSTM Layer}: Processes sequences in both forward and backward directions with 32 hidden units and L2 regularization (0.001), enabling the capture of long-range dependencies and contextual information \cite{kumar2023amc}
\item \textbf{Dropout Regularization}: Applied with a rate of 0.4 to prevent overfitting during training
\item \textbf{Global Average Pooling}: Aggregates sequence representations over time steps to create fixed-size feature vectors
\item \textbf{Dense Classification Layer}: A fully connected layer with 32 neurons and ReLU activation, followed by a final sigmoid output layer for binary sentiment classification \cite{pimpalkar2022mblstmglove}
\end{enumerate}

The model architecture is formally defined as:

\begin{equation}
h_t = \text{BiLSTM}(e_t, h_{t-1})
\end{equation}

\begin{equation}
\text{pooled} = \text{GlobalAvgPool}(h_1, h_2, ..., h_T)
\end{equation}

\begin{equation}
y = \sigma(W_2 \cdot \text{ReLU}(W_1 \cdot \text{pooled} + b_1) + b_2)
\end{equation}

where $e_t$ represents the GloVe embedding for token $t$, $h_t$ is the hidden state at time $t$, and $\sigma$ denotes the sigmoid activation function.


\subsubsection{Model Training Configuration}

The model training procedure follows established best practices for sentiment analysis tasks \cite{li2022glove}. The configuration includes:

\begin{itemize}
\item \textbf{Optimizer}: Adam optimizer with default learning rate (0.001)
\item \textbf{Loss Function}: Binary cross-entropy for binary sentiment classification
\item \textbf{Batch Size}: 64 samples per batch to balance computational efficiency and gradient stability
\item \textbf{Early Stopping}: Implemented with patience of 2 epochs monitoring validation loss to prevent overfitting
\item \textbf{Maximum Epochs}: 20 epochs with automatic termination via early stopping
\item \textbf{Regularization}: L2 regularization (0.001) applied to LSTM and dense layers, combined with dropout (0.4)
\end{itemize}

\subsection{Dataset Selection and Preprocessing}

\subsubsection{Yelp Polarity Reviews Dataset}

Our experiments utilize the Yelp Polarity Reviews dataset, a widely-recognized benchmark for binary sentiment classification tasks \cite{morris2020textattack}. The dataset characteristics are:

\begin{itemize}
\item \textbf{Training Set}: 560,000 reviews (280,000 positive, 280,000 negative)
\item \textbf{Test Set}: 38,000 reviews (19,000 positive, 19,000 negative)
\item \textbf{Class Distribution}: Balanced binary classification with stars 1-2 labeled as negative (class 0) and stars 4-5 labeled as positive (class 1)
\item \textbf{Domain}: Restaurant and business reviews from Yelp platform
\end{itemize}

\subsubsection{Text Preprocessing Pipeline}

The preprocessing pipeline ensures consistent input formatting and optimal model performance \cite{kumar2023amc}:

\begin{enumerate}
\item \textbf{Text Extraction}: Raw text extracted from TensorFlow Datasets format using UTF-8 decoding
\item \textbf{Tokenization}: Keras Tokenizer configured with vocabulary size limit of 50,000 most frequent words
\item \textbf{Sequence Conversion}: Text converted to integer sequences using fitted tokenizer vocabulary
\item \textbf{Padding and Truncation}: Sequences standardized to maximum length of 200 tokens using post-padding and post-truncation strategies
\item \textbf{Label Processing}: Sentiment labels converted to binary format (0 for negative, 1 for positive)
\end{enumerate}

\subsubsection{Justification for BiLSTM with GloVe Embeddings}

We selected the BiLSTM architecture with GloVe embeddings over transformer-based models like BERT for several compelling reasons. First, BiLSTM models offer significantly lower computational requirements while maintaining competitive performance for binary sentiment classification tasks \cite{li2022glove}\cite{kumar2023amc}. This efficiency is particularly valuable when developing robust defenses that require multiple training iterations. Second, BiLSTM's bidirectional processing capability effectively captures long-range dependencies in both directions, which is essential for sentiment analysis where emotional tone may depend on the entire sequence context \cite{pimpalkar2022mblstmglove}. Third, GloVe embeddings provide rich semantic representations through global word co-occurrence statistics, offering a strong foundation for identifying sentiment-bearing words \cite{li2022glove}. Finally, recent research has demonstrated that BiLSTM models with GloVe embeddings can achieve comparable or superior performance to transformer models on certain sentiment analysis tasks, particularly when dealing with adversarial perturbations \cite{kumar2023amc}\cite{pimpalkar2022mblstmglove}. This architecture thus represents an optimal balance between computational efficiency and model performance for our adversarial robustness investigation.



\subsection{TextAttack Framework Implementation}

\subsubsection{Framework Overview and Model Integration}

We implement adversarial attacks using the TextAttack framework, a comprehensive library for adversarial attacks in NLP that provides standardized implementations and evaluation metrics \cite{morris2020textattack}. The framework ensures reproducible and fair comparisons across different attack strategies through its modular design.

\subsubsection{Custom Model Wrapper Implementation}

To integrate our BiLSTM model with TextAttack, we implement a custom model wrapper that conforms to the framework's API requirements \cite{morris2020textattack}:

\begin{breakablealgorithm}{Custom Keras Model Wrapper for TextAttack Integration}
\begin{algorithmic}[1]
\Require BiLSTM model, tokenizer, max\_length=200 \Comment{Target model, text processor, sequence length}
\Ensure Compatible model wrapper for TextAttack \Comment{Returns wrapper instance}

\State Initialize wrapper with model, tokenizer, and max\_length parameters
\Function{\_\_call\_\_}{text\_input\_list} \Comment{Main processing function}
    \State $sequences \gets tokenizer.texts\_to\_sequences(text\_input\_list)$ \Comment{Convert text to token IDs}
    \State $padded \gets pad\_sequences(sequences, maxlen=max\_length, padding='post')$ \Comment{Pad/truncate sequences}
    \State $preds \gets model.predict(padded)$ \Comment{Get model predictions}
    \State \Return $concatenate([1 - preds, preds], axis=1)$ \Comment{Convert to binary class probabilities}
\EndFunction
\end{algorithmic}
\end{breakablealgorithm}

\subsection{Adversarial Attack Implementation}

\subsubsection{TextFooler Attack Configuration}

TextFooler represents a sophisticated synonym-based substitution attack that leverages semantic similarity constraints to generate natural adversarial examples \cite{jin2019textfooler}. Our implementation configures the attack with the following components:

\begin{itemize}
\item \textbf{Transformation}: WordSwapEmbedding with maximum of 50 candidate substitutions per word
\item \textbf{Constraints}:
\begin{itemize}
\item RepeatModification to prevent redundant transformations
\item StopwordModification to preserve linguistic function words
\item WordEmbeddingDistance with minimum cosine similarity of 0.5
\item PartOfSpeech constraint allowing verb-noun swaps
\item UniversalSentenceEncoder with threshold of 0.840845057 and angular metric
\end{itemize}
\item \textbf{Search Method}: GreedyWordSwapWIR with deletion-based word importance ranking
\item \textbf{Goal Function}: UntargetedClassification for general misclassification
\end{itemize}

\subsubsection{TextBugger Attack Configuration}

TextBugger employs a multi-level transformation approach combining character-level and word-level modifications \cite{li2018textbugger}. The attack implementation includes:

\begin{itemize}
\item \textbf{Character-Level Transformations}:
\begin{itemize}
\item WordSwapRandomCharacterInsertion (space insertion)
\item WordSwapRandomCharacterDeletion (character removal)
\item WordSwapNeighboringCharacterSwap (adjacent character swapping)
\item WordSwapHomoglyphSwap (visually similar character substitution)
\end{itemize}
\item \textbf{Word-Level Transformation}: WordSwapEmbedding with top-5 nearest neighbors
\item \textbf{Constraints}: UniversalSentenceEncoder with threshold of 0.8 to maintain semantic similarity
\item \textbf{Search Strategy}: GreedyWordSwapWIR for efficient perturbation identification
\end{itemize}

\subsubsection{DeepWordBug Attack Configuration}

DeepWordBug operates as a black-box attack using character-level perturbations with strategic scoring \cite{gao2018deepwordbug}. The implementation specifications include:

\begin{itemize}
\item \textbf{Transformation Set}:
\begin{itemize}
\item WordSwapNeighboringCharacterSwap for adjacent character exchanges
\item WordSwapRandomCharacterSubstitution for single character replacement
\item WordSwapRandomCharacterDeletion for character removal
\item WordSwapRandomCharacterInsertion for character addition
\end{itemize}
\item \textbf{Constraints}:
\begin{itemize}
\item LevenshteinEditDistance with maximum of 30 character modifications
\item StopwordModification to preserve functional words
\item RepeatModification to avoid duplicate transformations
\end{itemize}
\item \textbf{Search Method}: GreedyWordSwapWIR for importance-based token ranking
\end{itemize}

\begin{table}[H]
\caption{Comparison of Implemented Adversarial Attack Methods}
\label{tab:attack_comparison}
\centering
\footnotesize
\begin{tabular}{@{\hspace{0.3cm}}p{2.5cm}p{3.5cm}p{3.5cm}p{3.5cm}@{\hspace{0.3cm}}}
\toprule
\textbf{Feature} & \textbf{TextFooler} & \textbf{TextBugger} & \textbf{DeepWordBug} \\
\midrule
Attack Level & Word-level & Word and character-level & Character-level \\
\midrule
Transformation & Word substitution with semantically similar alternatives & Character insertion, deletion, swapping, and word substitution & Character swapping, substitution, deletion, insertion \\
\midrule
Key Constraint & Universal Sentence Encoder (0.84 threshold) & Universal Sentence Encoder (0.8 threshold) & Levenshtein Edit Distance (max 30) \\
\midrule
Search Method & Greedy Word Importance Ranking & Greedy Word Importance Ranking & Greedy Word Importance Ranking \\
\midrule
Attack Type & Black-box & Black-box & Black-box \\
\midrule
Primary Strength & Preserves semantic meaning while maximizing misclassification & Effective against real-world systems with multiple transformation types & Computationally efficient with minimal model access requirements \\
\bottomrule
\end{tabular}
\end{table}


\subsubsection{Rationale for Selected Attack Methods}

The three attack methods implemented in this study—TextFooler, TextBugger, and DeepWordBug—were strategically selected to represent diverse adversarial approaches against sentiment analysis models \cite{jin2019textfooler}\cite{li2018textbugger}\cite{gao2018deepwordbug}. These methods collectively cover the spectrum of text perturbation strategies: TextFooler operates primarily at the word level with semantic constraints, TextBugger combines both character and word-level modifications, and DeepWordBug focuses exclusively on character-level perturbations \cite{morris2020textattack}. This diversity enables comprehensive evaluation of model vulnerabilities across different linguistic levels. Additionally, all three attacks have demonstrated high success rates against state-of-the-art NLP models while maintaining text fluency and semantic preservation \cite{jin2019textfooler}\cite{li2018textbugger}. Furthermore, these methods represent practical black-box attack scenarios that reflect real-world threats, where adversaries typically have limited access to model internals \cite{morris2020textattack}\cite{gao2018deepwordbug}. By implementing these complementary attack strategies through the standardized TextAttack framework, we ensure fair comparison and reproducible evaluation of adversarial vulnerabilities in sentiment analysis models.


\subsection{Virtual Adversarial Training Defense Implementation}

\subsubsection{VAT Theoretical Foundation}

Virtual Adversarial Training enhances model robustness by generating virtual adversarial examples during training through gradient-based perturbation generation \cite{miyato2017virtual}. A key advantage of VAT is that it does not require labeled adversarial examples, as the definition of virtual adversarial perturbation only requires input $x$ and does not require label $y$, making it applicable to semi-supervised learning scenarios \cite{miyato2017virtual}.

VAT operates by identifying directions in the input space where the model's output distribution is most sensitive, then regularizing the model to be smooth in these directions \cite{miyato2017virtual}. The method generates perturbations that maximize the KL divergence between the original and perturbed predictions, effectively finding the locally adversarial direction for each input \cite{miyato2017virtual}. This approach allows VAT to improve model robustness without requiring explicit adversarial examples, making it computationally efficient compared to traditional adversarial training methods \cite{miyato2017virtual}.


\subsubsection{VAT Model Architecture and Implementation}

Our VAT implementation extends the base BiLSTM model with adversarial regularization capabilities by applying perturbations at the embedding level rather than directly to the input tokens:

\begin{breakablealgorithm}{Virtual Adversarial Training (VAT) - Verified Implementation}
\begin{algorithmic}[1]
\Require Base model $f_\theta$, perturbation magnitude $\epsilon=2.0$, step size $\xi=1e-6$, iterations $K=1$ 
\Ensure Robust model $f_\theta$ with adversarial regularization

\State Initialize model parameters $\theta$ from base model
\Function{TrainingStep}{input batch $(X, y)$}
    \State \textbf{Standard Forward Pass:}
    \State $y_{\text{pred}} \gets f_\theta(X)$ \Comment{Compute model predictions}
    \State $\mathcal{L}_{\text{CE}} \gets \text{CrossEntropy}(y_{\text{pred}}, y)$ \Comment{Classification loss}
    
    \State \textbf{Virtual Adversarial Direction Finding:} 
    \State $d \gets \mathcal{N}(0, I)$ \Comment{Initialize random unit vector}
    \For{$k \gets 1$ to $K$} \Comment{Power iteration for adversarial direction}
        \State $d \gets \xi \cdot \frac{d}{\|d\|_2}$ \Comment{Normalize and scale perturbation}
        \State $\tilde{X} \gets X + d$ \Comment{Apply temporary perturbation}
        \State $\mathcal{L}_{\text{KL}} \gets \text{KL}(f_\theta(X) \| f_\theta(\tilde{X}))$ \Comment{KL divergence}
        \State $d \gets \nabla_d \mathcal{L}_{\text{KL}}$ \Comment{Compute gradient w.r.t. perturbation}
        \State $d \gets \text{stop\_gradient}(\frac{d}{\|d\|_2})$ \Comment{Prevent gradient flow + normalize}
    \EndFor
    
    \State \textbf{Final Virtual Adversarial Perturbation:}
    \State $r_{\text{adv}} \gets \epsilon \cdot \frac{d}{\|d\|_2}$ \Comment{Scale by $\epsilon$}
    
    \State \textbf{VAT Regularization Loss:}
    \State $\mathcal{L}_{\text{VAT}} \gets \text{KL}(f_\theta(X) \| f_\theta(X + r_{\text{adv}}))$ \Comment{Virtual adversarial loss}
    
    \State \textbf{Combined Loss Optimization:}
    \State $\theta \gets \theta - \eta \nabla_\theta (\mathcal{L}_{\text{CE}} + \alpha \mathcal{L}_{\text{VAT}})$ \Comment{$\alpha=1$}
    
    \State \Return Updated parameters $\theta$
\EndFunction
\end{algorithmic}
\end{breakablealgorithm}


\subsubsection{VAT Hyperparameter Configuration}

The VAT implementation employs carefully selected hyperparameters based on the original VAT methodology and empirical optimization for text classification tasks \cite{miyato2017virtual}:

\begin{itemize}
\item \textbf{Perturbation Magnitude ($\epsilon$)}: 2.0, selected through empirical optimization for sentiment analysis tasks. While the original VAT paper notes that optimal $\epsilon$ values are task-dependent, our choice reflects the embedding space scale of GloVe vectors and the need for sufficient perturbation magnitude to achieve regularization effects \cite{miyato2017virtual}
\item \textbf{Step Size ($\xi$)}: 1e-6, chosen to ensure numerical stability during power iteration while maintaining sufficient gradient signal for perturbation direction finding \cite{miyato2017virtual}
\item \textbf{Power Iterations ($ip$)}: 1 iteration for computational efficiency. The original VAT framework demonstrates that even a single power iteration can provide effective approximation of the adversarial direction when combined with appropriate step size selection \cite{miyato2017virtual}
\item \textbf{VAT Loss Weight ($\alpha$)}: 1.0, following the original VAT recommendation where superior performance is achieved by tuning only $\epsilon$ while fixing $\alpha=1$ \cite{miyato2017virtual}
\item \textbf{Gradient Clipping}: [-0.5, 0.5] range for training stability during adversarial training
\item \textbf{Learning Rate Schedule}: Cosine decay from 1e-4 over 1000 steps for stable convergence
\end{itemize}

The original VAT paper emphasizes that there are only two scalar-valued hyperparameters requiring tuning: the norm constraint $\epsilon > 0$ for the adversarial direction and the regularization coefficient $\alpha > 0$ \cite{miyato2017virtual}. Our hyperparameter selection aligns with this principle, focusing primarily on $\epsilon$ optimization while maintaining the recommended $\alpha=1$ setting.


\subsubsection{Justification for Virtual Adversarial Training}

We selected Virtual Adversarial Training (VAT) as our defense mechanism over alternatives like traditional adversarial training for several key advantages in the NLP domain \cite{miyato2017virtual}. First, VAT generates adversarial perturbations without requiring labeled examples, making it applicable to semi-supervised learning scenarios where labeled data may be limited \cite{miyato2017virtual}\cite{morris2020textattack}. This is particularly valuable for sentiment analysis tasks where obtaining large quantities of labeled data can be resource-intensive. Second, VAT creates perturbations on-the-fly using gradient information, offering computational efficiency compared to methods that require explicit adversarial example generation \cite{miyato2017virtual}. Third, VAT has demonstrated superior performance in regularizing neural networks by smoothing the model's decision boundaries in the direction most vulnerable to adversarial perturbations \cite{miyato2017virtual}. Finally, unlike traditional adversarial training that focuses on specific attack types, VAT improves general robustness by considering the worst-case perturbations within a small neighborhood of each input \cite{miyato2017virtual}\cite{morris2020textattack}. This approach aligns with our goal of developing sentiment analysis models that are robust against diverse and potentially unknown adversarial attacks.

\subsection{Character-Level Purification Defense Implementation}

\subsubsection{Purification Framework Design}

To address the limitations of VAT against character-level attacks, we implement a lightweight character-level purification defense that operates as a preprocessing layer before text reaches the VAT-trained model. The purification framework targets the specific character-level transformations employed by TextBugger and DeepWordBug attacks, providing complementary protection to VAT's embedding-space regularization.

\subsubsection{Purification Algorithm Components}

Our character purification implementation consists of four main components designed to counter specific attack strategies:

\begin{enumerate}
\item \textbf{Character Substitution Correction}: Reverses common homoglyph substitutions (0→o, 1→l, @→a, \$→s) frequently used in character-level attacks
\item \textbf{Repetition Normalization}: Removes excessive character repetitions while preserving legitimate doubles (ll, ss, etc.) using regex pattern matching
\item \textbf{Single Edit Correction}: Attempts single character deletion, insertion, and substitution to restore out-of-vocabulary words to valid vocabulary entries
\item \textbf{Adjacent Swap Correction}: Fixes character transpositions commonly used in swap-based attacks through systematic adjacent character exchange
\end{enumerate}

\begin{breakablealgorithm}{Character-Level Purification Algorithm}
\begin{algorithmic}[1]
\Require Input text $T$, vocabulary $V$, character substitution map $S$
\Ensure Purified text $T'$ with character-level perturbations corrected

\Function{PurifyText}{$T$}
    \State $words \gets \text{split}(T, \text{' '})$ \Comment{Tokenize input text}
    \State $purified\_words \gets []$
    
    \For{each $word$ in $words$}
        \If{$word.lower() \in \text{protected\_stopwords}$}
            \State $purified\_words.\text{append}(word)$ \Comment{Skip protected words}
            \State \textbf{continue}
        \EndIf
        
        \If{$word.lower() \in V$}
            \State $purified\_words.\text{append}(word)$ \Comment{Valid vocabulary word}
            \State \textbf{continue}
        \EndIf
        
        \State $corrected \gets \text{ApplyCharacterCorrections}(word, S, V)$
        \State $purified\_words.\text{append}(corrected)$
    \EndFor
    
    \State \Return $\text{join}(purified\_words, \text{' '})$
\EndFunction

\Function{ApplyCharacterCorrections}{$word$, $S$, $V$}
    \State \textbf{Step 1:} Apply homoglyph substitutions from map $S$
    \State \textbf{Step 2:} Remove excessive character repetitions
    \State \textbf{Step 3:} Try single character deletion/insertion/substitution
    \State \textbf{Step 4:} Try adjacent character swaps
    \State \Return $corrected\_word$ if valid, else $original\_word$
\EndFunction
\end{algorithmic}
\end{breakablealgorithm}

\subsubsection{Implementation Specifications}

The purification defense employs the following configuration optimized for computational efficiency:

\begin{itemize}
\item \textbf{Vocabulary Reference}: Uses the same tokenizer vocabulary as the target model (50,000 words) for consistency
\item \textbf{Protected Words}: Preserves stopwords and function words that are typically protected in attacks, matching the TextAttack framework's stopword list

\item \textbf{Character Frequency Analysis}: Prioritizes corrections using character frequency analysis from vocabulary sample (50000 words)
\item \textbf{Processing Strategy}: Sequential application of correction methods with early termination upon successful vocabulary match
\end{itemize}

\subsubsection{Integration with VAT Defense}

The hybrid defense architecture combines character-level purification with VAT through a two-stage process:

\begin{enumerate}
\item \textbf{Stage 1 - Input Purification}: Character-level corrections applied to adversarial text before model processing
\item \textbf{Stage 2 - VAT Robustness}: Purified text processed by VAT-trained model with embedding-space regularization
\end{enumerate}

This layered approach addresses vulnerabilities at multiple linguistic levels: purification handles discrete character-level perturbations while VAT provides continuous embedding-space robustness.


\subsection{Evaluation Metrics and Experimental Design}

\subsubsection{Performance Metrics}

Our evaluation framework employs standard metrics for adversarial robustness assessment \cite{morris2020textattack}:

\begin{itemize}
\item \textbf{Original Accuracy}: Model performance on clean, unperturbed test data
\item \textbf{Accuracy Under Attack (AuA)}: Model performance on adversarial examples
\item \textbf{Attack Success Rate (ASR)}: Percentage of successful adversarial examples
\item \textbf{Accuracy Recovery}: Performance improvement achieved through defense mechanisms
\end{itemize}

The metrics are formally defined as:

\begin{equation}
\text{ASR} = \frac{\text{Number of Successful Attacks}}{\text{Total Attack Attempts}} \times 100%
\end{equation}

\begin{equation}
\text{AuA} = \frac{\text{Correctly Classified Adversarial Examples}}{\text{Total Adversarial Examples}} \times 100%
\end{equation}

\begin{equation}
\text{Defense Effectiveness} = \text{AuA}{defended} - \text{AuA}{original}
\end{equation}

\subsubsection{Experimental Protocol}

Our experimental design ensures reproducible and statistically significant results:

\begin{enumerate}
\item \textbf{Sample Selection}: Evaluation performed on samples correctly classified by the original model to measure attack effectiveness
\item \textbf{Sample Sizes}: 1000 samples initially selected for each attack method, with evaluation performed on correctly classified samples (resulting in approximately 957 valid samples after filtering)

\item \textbf{Attack Constraints}: All attacks limited to preserve semantic meaning and grammatical correctness
\item \textbf{Defense Training}: VAT models trained for 8 epochs using Adam optimizer with cosine decay learning rate schedule (initial: 1e-4), batch size of 32, on 5,000 sample subset with 20\% validation split

\end{enumerate}

\subsection{Implementation Environment}

\subsubsection{Software Dependencies and Versions}

The experimental environment ensures reproducibility through standardized software configurations:

\begin{itemize}
\item \textbf{Python}: 3.8+ with TensorFlow 2.x framework
\item \textbf{TextAttack}: Latest version for attack implementations
\item \textbf{TensorFlow Datasets}: For Yelp Polarity Reviews data loading
\item \textbf{NumPy and SciPy}: For numerical computations and statistical analysis
\item \textbf{Transformers}: Version 4.41.0 for compatibility with TextAttack
\end{itemize}


This comprehensive methodology provides a robust framework for evaluating adversarial vulnerabilities in sentiment analysis models and assessing the effectiveness of Virtual Adversarial Training as a defense mechanism. The systematic approach ensures fair comparisons across different attack strategies while maintaining scientific rigor in experimental design and evaluation.

\section{Experimental Results}

This section presents detailed metrics for three adversarial attacks—TextFooler, TextBugger, and DeepWordBug—on our BiLSTM–GloVe sentiment model, both before and after applying Virtual Adversarial Training (VAT) defense \cite{jin2019textfooler}\cite{li2018textbugger}\cite{gao2018deepwordbug}.

\subsection{TextFooler Results}
\begin{table}[H]
\caption{TextFooler Attack and Hybrid Defense Metrics}
\label{tab:textfooler_results}
\centering
\footnotesize
\begin{tabular}{@{\hspace{0.3cm}}l*{3}{S[table-format=3.2]}@{\hspace{0.3cm}}}
\toprule
\textbf{Metric} & {\textbf{Original}} & {\textbf{VAT Only}} & {\textbf{VAT+Purif}} \\
\midrule
Original Accuracy (\%) & 95.70 & 95.70 & 95.70 \\
Accuracy Under Attack (AuA) (\%) & 3.97 & 49.74 & 49.74 \\
Attack Success Rate (ASR) (\%) & 96.03 & 50.26 & 50.26 \\
Successful Attacks & 919.00 & 481.00 & 481.00 \\
Failed Attacks & 38.00 & 476.00 & 476.00 \\
Total Attempted & 957.00 & 957.00 & 957.00 \\
ASR Reduction (\%) & {---} & 45.77 & 45.77 \\
\bottomrule
\end{tabular}
\end{table}



\FloatBarrier

TextFooler's synonym-based substitutions achieved a 96.03\% ASR on the original model. VAT defense improved accuracy under attack to 49.74\%, reducing ASR by 45.77\%. The purification component showed no additional improvement for word-level attacks, as expected since character-level corrections do not address semantic substitutions \cite{jin2019textfooler}.


\subsection{TextBugger Results}
\begin{table}[H]
\caption{TextBugger Attack and Hybrid Defense Metrics}
\label{tab:textbugger_results}
\centering
\footnotesize
\begin{tabular}{@{\hspace{0.3cm}}l*{3}{S[table-format=3.2]}@{\hspace{0.3cm}}}
\toprule
\textbf{Metric} & {\textbf{Original}} & {\textbf{VAT Only}} & {\textbf{VAT+Purif}} \\
\midrule
Original Accuracy (\%) & 95.70 & 95.70 & 95.70 \\
Accuracy Under Attack (AuA) (\%) & 7.31 & 52.98 & 57.68 \\
Attack Success Rate (ASR) (\%) & 92.69 & 47.02 & 42.32 \\
Successful Attacks & 887.00 & 450.00 & 405.00 \\
Failed Attacks & 70.00 & 507.00 & 552.00 \\
Total Attempted & 957.00 & 957.00 & 957.00 \\
ASR Reduction (\%) & {---} & 45.66 & 50.37 \\
\bottomrule
\end{tabular}
\end{table}

\FloatBarrier

TextBugger achieved a 92.69\% ASR through mixed character and word-level modifications. VAT defense improved accuracy to 52.98\% (45.66\% ASR reduction), while the hybrid VAT+Purification approach achieved 57.68\% accuracy with 50.37\% ASR reduction. The additional 4.71\% improvement demonstrates purification's effectiveness against the character-level components of TextBugger \cite{li2018textbugger}.


\subsection{DeepWordBug Results}
\begin{table}[H]
\caption{DeepWordBug Attack and Hybrid Defense Metrics}
\label{tab:deepwordbug_results}
\centering
\footnotesize
\begin{tabular}{@{\hspace{0.3cm}}l*{3}{S[table-format=3.2]}@{\hspace{0.3cm}}}
\toprule
\textbf{Metric} & {\textbf{Original}} & {\textbf{VAT Only}} & {\textbf{VAT+Purif}} \\
\midrule
Original Accuracy (\%) & 95.70 & 95.70 & 95.70 \\
Accuracy Under Attack (AuA) (\%) & 7.42 & 43.68 & 84.33 \\
Attack Success Rate (ASR) (\%) & 92.58 & 56.32 & 15.67 \\
Successful Attacks & 886.00 & 539.00 & 150.00 \\
Failed Attacks & 71.00 & 418.00 & 807.00 \\
Total Attempted & 957.00 & 957.00 & 957.00 \\
ASR Reduction (\%) & {---} & 36.26 & 76.91 \\
\bottomrule
\end{tabular}
\end{table}

\FloatBarrier

DeepWordBug achieved a 92.58\% ASR through character-level perturbations. VAT defense improved accuracy to 43.68\% (36.26\% ASR reduction), while the hybrid approach dramatically improved accuracy to 84.33\% with 76.91\% ASR reduction. This substantial 40.65\% additional improvement validates the effectiveness of character-level purification against pure character-level attacks \cite{gao2018deepwordbug}.


\subsection{Overall ASR Reduction Comparison}
\begin{table}[H]
\caption{Comprehensive Defense Effectiveness Comparison}
\label{tab:asr_comparison}
\centering
\footnotesize
\begin{tabular}{@{\hspace{0.3cm}}l*{4}{S[table-format=2.2]}@{\hspace{0.3cm}}}
\toprule
\textbf{Attack Method} & {\textbf{Original ASR (\%)}} & {\textbf{VAT ASR (\%)}} & {\textbf{VAT+Purif ASR (\%)}} & {\textbf{Total Reduction (\%)}} \\
\midrule
TextFooler & 96.03 & 50.26 & 50.26 & 45.77 \\
TextBugger & 92.69 & 47.02 & 42.32 & 50.37 \\
DeepWordBug & 92.58 & 56.32 & 15.67 & 76.91 \\
\bottomrule
\end{tabular}
\end{table}
\FloatBarrier


\subsection{Hybrid Defense Analysis}

\subsubsection{Defense Effectiveness by Attack Granularity}

The experimental results demonstrate a clear hierarchy in defense effectiveness that validates our theoretical expectations:

\begin{itemize}
\item \textbf{Word-level attacks (TextFooler)}: Purification provides no additional benefit (0\% improvement) since character-level corrections cannot address semantic word substitutions
\item \textbf{Hybrid attacks (TextBugger)}: Moderate purification improvement (4.71\% additional ASR reduction) due to character-level components being corrected while word-level components remain challenging
\item \textbf{Character-level attacks (DeepWordBug)}: Substantial purification improvement (40.65\% additional ASR reduction) validating the targeted approach against pure character-level perturbations
\end{itemize}

\subsubsection{Synergistic Defense Effects}

The combination of VAT and character purification creates a comprehensive defense strategy where:

\begin{enumerate}
\item \textbf{Complementary Coverage}: VAT provides embedding-space regularization while purification handles input-space character-level perturbations
\item \textbf{Layered Protection}: Two-stage defense process addresses vulnerabilities at multiple linguistic levels
\item \textbf{Attack-Specific Effectiveness}: Defense performance correlates with attack granularity, achieving highest effectiveness against character-level attacks (76.91\% total ASR reduction)
\end{enumerate}

The results confirm that character-level purification is most effective against attacks that operate at the character level, providing minimal benefit for semantic word-level attacks but substantial improvement for character-level perturbations.



\end{document}